\documentclass[10pt,a4paper]{amsart}
\usepackage[margin=19mm]{geometry}
\usepackage{amsmath,amssymb,mathtools}
\usepackage[scr=rsfs,cal=euler]{mathalfa}
\usepackage{xfrac}
\usepackage[unicode,hidelinks,bookmarks=false]{hyperref}
\newcommand{\ie}{i.e.\ }
\newcommand{\cf}{cf.\ }
\newcommand{\resp}{respectively\ }
\newcommand{\Subsection}{\subsection}
\providecommand{\nobreakdash}{\nobreak\hbox{-}}
\setlength{\emergencystretch}{2em}
\multlinegap0pt
\usepackage{amssymb}
\usepackage{bbm}
\usepackage[shortlabels]{enumitem}
\usepackage{mathtools}
\usepackage{caption}
\newtheorem{lemma}{Lemma}
\newcommand{\lip}{\mathrm{Lip}}
\title{Existence of spectral submanifolds in time delay systems}
\author{Gergely Buza, George Haller}
\date{Source: arXiv:2512.22062v1 (CC BY 4.0)}
\begin{document}
\maketitle

\noindent\emph{Source and licence.} Existence of spectral submanifolds in time delay systems. Version arXiv:2512.22062v1 (CC BY 4.0), distributed by arXiv under the Creative Commons Attribution 4.0 International licence, http://creativecommons.org/licenses/by/4.0/, as recorded in the arXiv licence metadata for this exact version and shown on the abstract page https://arxiv.org/abs/2512.22062v1. Full licence notice: \texttt{licenses/arxiv-2512.22062-CC-BY-4.0.txt}.\par\smallskip

\noindent\textit{Editorial scope.} Counted unit: the article's lemma lemma:nonlinearity\_sf (source0.tex lines 1792--1803). The article's complete original proof of it is delivered with it (source0.tex lines 1805--1824, one \texttt{proof} environment of 20 lines). No mathematics is written, repaired or re-derived: the statement and the proof are the article's own lines, in the article's own notation, and no retained line is substituted or re-flowed.  Retained uncounted as the source-local context the proof needs: lines 1688--1691; lines 1712--1716; lines 1743--1746; lines 1754--1756.  Nothing was omitted from any retained span, so the package carries no omission record.  No retained line contains a citation, so the wrapper carries no bibliography and no bibliography entry.  No retained line contains a figure environment, an image-inclusion command or drawing code, so no image is delivered and the package declares no asset dependency.  Editorial formatting only, in addition: an AMS \texttt{amsart} preamble that restates the article's own declarations and macros used by the retained lines, namely the environments \texttt{lemma} on separate counters, together with the standard packages those lines need.  The retained environments print in this document as Lemma 1 in order of appearance, whereas the article prints the same items with the numbering of its own full document.  The complete native source of the article as distributed in the arXiv e-print for this exact version is bundled unchanged as \texttt{sources/191712/source0.tex}.
\begin{equation}
    R_\rho (u) : =  \chi_\rho ( |P_\Sigma u |_X) \chi_\rho (|P_{\Sigma'} u |_X) R(u), \qquad u \in O,
    \label{eq:Rrho}
\end{equation}

\begin{lemma} \label{lemma:Lipschitz}
    Let $R_\rho :X \to X^{\odot *}$ be as in \eqref{eq:Rrho}, extended to the whole of $X$ as above.
    Denote by $\lip_\delta(R)$ the Lipschitz constant of $R$ on the (closed) ball $\overline{B_\delta(0)}^X$ of radius $\delta$ centered at the origin, and by $\lip(R_\rho)$ the global Lipschitz constant of $R_\rho:X \to X^{\odot *} $.
    Then, there exist  constants $C,c >0 $ such that $\lip(R_\rho) \leq  C \, \lip_{c \rho}(R) = o(1)$ as $\rho \to 0$.\footnote{A straightforward computation shows that we may take $c = 4$ and $C = 2(1+2\Vert P_\Sigma \Vert)$.}
\end{lemma}

\begin{equation}
    \varphi_t^\rho(u) = T(t) u + \imath^{-1} \int_0^t T^{\odot *} (t-s) R_\rho \circ \varphi_s^\rho(u) \, ds,  \qquad (t,u) \in \mathcal{D}^{\varphi^\rho}.
    \label{eq:varphirho}
\end{equation}

\begin{lemma}[Global existence for $\varphi^\rho$] \label{lemma:globalexist}
    Equation \eqref{eq:varphirho} defines a maximal semiflow $\varphi^\rho$ with domain $\mathcal{D}^{\varphi^\rho} = \mathbb{R}^{\geq 0} \times X$.
\end{lemma}

\begin{lemma} \label{lemma:nonlinearity_sf}
    Let $\omega > \omega(A)$, $\omega \neq 0$, and $M>0$ be as in the proof of Lemma~\ref{lemma:globalexist}, and fix $t \geq 0$.
    Then $\varphi_t^\rho : X \to X$ is globally Lipschitz for $\rho > 0$ sufficiently small (depending on $t$), with
    \begin{displaymath}
       \sup_{s \in [0,t]} \lip\big( \varphi_s^\rho\big) \leq 2M e^{t \omega}.
    \end{displaymath}
    Moreover,
    \begin{equation}
         \lip\big(N_\rho(t,\cdot)\big) \leq \frac{2M^2}{\omega} \left( e^{\omega t} - 1 \right)  e^{\omega t} \lip(R_\rho).
         \label{eq:LipN}
    \end{equation}
\end{lemma}

\begin{proof}
    The case $t =0$ is trivial, so assume $t > 0$.
    Then, for $u,v \in X$,
    \begin{displaymath}
        |\varphi_t^\rho (u) - \varphi_t^\rho(v)|_X \leq Me^{\omega t}  |u-v|_X + \frac{M}{\omega} \left( e^{\omega t} - 1 \right) \lip(R_\rho) \sup_{s \in [0,t]}  |\varphi_s^\rho (u) - \varphi_s^\rho(v)|_X.
    \end{displaymath}
    Taking the supremum over $t \in [0,T]$ and choosing $\rho$ small enough such that (c.f.\ Lemma~\ref{lemma:Lipschitz})
    \begin{equation}
        \lip(R_\rho) < \frac{\omega}{ 2 M \left( e^{\omega t} - 1 \right)},
        \label{eq:replaced}
    \end{equation}
    we obtain 
    \begin{displaymath}
        \sup_{t \in [0,T]} |\varphi_t^\rho (u) - \varphi_t^\rho(v)|_X \leq 2Me^{\omega T}  |u-v|_X .
    \end{displaymath}
    Moreover, 
    \begin{displaymath}
        |N_\rho(t,u ) - N_\rho(t,v)|_X \leq \frac{M}{\omega} \left( e^{\omega t} - 1 \right) \lip(R_\rho) 2Me^{\omega t}  |u-v|_X. \qedhere
    \end{displaymath}
\end{proof}

\end{document}
